\documentclass[11pt,a4paper]{article}

\usepackage[utf8]{inputenc}
\usepackage[T1]{fontenc}
\usepackage{geometry}
\usepackage{booktabs}
\usepackage{xcolor}
\usepackage{colortbl}
\usepackage{fancyhdr}
\usepackage{titlesec}
\usepackage{hyperref}
\usepackage{tcolorbox}
\usepackage{enumitem}
\usepackage{array}
\usepackage{parskip}
\usepackage{amsmath}
\usepackage{amssymb}

\geometry{left=3.0cm,right=3.0cm,top=3.0cm,bottom=2.8cm}

\definecolor{ink}{RGB}{18,18,18}
\definecolor{primary}{RGB}{10,50,100}
\definecolor{rule}{RGB}{10,50,100}
\definecolor{muted}{RGB}{100,108,120}
\definecolor{light}{RGB}{246,248,251}
\definecolor{accent}{RGB}{180,20,30}
\definecolor{rowalt}{RGB}{240,244,250}
\definecolor{greenlight}{RGB}{240,253,244}
\definecolor{greenborder}{RGB}{22,101,52}

\hypersetup{colorlinks=true,linkcolor=primary,urlcolor=primary,pdfborder={0 0 0}}

\titleformat{\section}
  {\large\bfseries\color{primary}}{}{0em}{}
  [\vspace{2pt}{\color{rule}\rule{\linewidth}{0.5pt}}]
\titlespacing{\section}{0pt}{1.6em}{0.6em}

\titleformat{\subsection}
  {\normalsize\bfseries\color{ink}}{}{0em}{}
\titlespacing{\subsection}{0pt}{0.9em}{0.2em}

\pagestyle{fancy}
\fancyhf{}
\fancyhead[L]{\small\color{muted}Confidencial}
\fancyhead[R]{\small\color{muted}Mantenimiento Predictivo de Flota --- Fundamentos}
\fancyfoot[C]{\small\color{muted}\thepage}
\renewcommand{\headrulewidth}{0pt}

\tcbuselibrary{skins,breakable}

\newtcolorbox{concepto}{
  colback=light, colframe=primary, boxrule=0.6pt,
  arc=2pt, left=8pt, right=8pt, top=5pt, bottom=5pt,
  breakable, fontupper=\small}

\newtcolorbox{resultado}{
  colback=white, colframe=primary, boxrule=0pt,
  leftrule=3pt, arc=0pt,
  left=8pt, right=4pt, top=4pt, bottom=4pt,
  breakable, fontupper=\small\color{ink}}

\newtcolorbox{critico}{
  colback=white, colframe=accent, boxrule=0pt,
  leftrule=3pt, arc=0pt,
  left=8pt, right=4pt, top=4pt, bottom=4pt,
  breakable, fontupper=\small\color{ink}}

\newtcolorbox{ahorro}{
  colback=greenlight, colframe=greenborder, boxrule=0.6pt,
  arc=2pt, left=8pt, right=8pt, top=5pt, bottom=5pt,
  breakable, fontupper=\small\color{ink}}

\setlength{\parskip}{0.45em}
\setlength{\parindent}{0em}
\renewcommand{\arraystretch}{1.3}

\setlist[itemize]{nosep, topsep=3pt, leftmargin=1.5em,
  label={\color{primary}\textbullet}}
\setlist[enumerate]{nosep, topsep=3pt, leftmargin=1.8em}

\begin{document}

%--- PORTADA ------------------------------------------------------------------
\thispagestyle{empty}

\vspace*{1.4cm}

{\color{muted}\small CONFIDENCIAL \quad|\quad Junio 2026}

\vspace{0.7cm}

{\fontsize{28}{33}\selectfont\bfseries\color{primary}
Mantenimiento Predictivo\\de Flota Vehicular}

\vspace{0.4cm}

{\fontsize{14}{18}\selectfont\color{ink}
Fundamentos matematicos, fisica de falla\\
y arquitectura de estimacion online}

\vspace{0.3cm}

{\color{rule}\rule{\linewidth}{0.8pt}}

\vspace{0.5cm}

{\color{muted}\normalsize
Este documento describe el marco cientifico del modulo de mantenimiento
que estamos construyendo: que matematicas usamos, por que son las correctas,
en que condiciones de datos opera el sistema desde el primer dia,
y que ahorro concreto y cuantificable produce su adopcion.
No es una propuesta comercial. Es la fisica del problema.}

\vspace{2.2cm}

\begin{tabular}{@{}p{0.55\linewidth}p{0.4\linewidth}@{}}
\small\color{muted}\textbf{I.} \color{ink} El problema y la asimetria de costos
  & \small\color{muted}\textbf{VI.} \color{ink} Estimacion online en regimen de streaming \\[3pt]
\small\color{muted}\textbf{II.} \color{ink} Confiabilidad y vida util remanente
  & \small\color{muted}\textbf{VII.} \color{ink} Marco jerarquico bayesiano \\[3pt]
\small\color{muted}\textbf{III.} \color{ink} Fisica de falla por componente
  & \small\color{muted}\textbf{VIII.} \color{ink} Cuantificacion del ahorro \\[3pt]
\small\color{muted}\textbf{IV.} \color{ink} Inferencia con censura y covariables
  & \small\color{muted}\textbf{IX.} \color{ink} Por que este enfoque domina \\[3pt]
\small\color{muted}\textbf{V.} \color{ink} Procesos de degradacion y RUL
  & \\[3pt]
\end{tabular}

\vspace{2.5cm}

\begin{concepto}
\textbf{El hilo conductor.}
Una pieza es un sistema fisico cuyo estado se degrada; la falla es el cruce
de un umbral. Todo lo demas ---estadistica de vida, fisica de desgaste,
decision optima--- son tres lentes sobre ese mismo hecho. Las tres lentes
convergen en un solo motor de inferencia que corre en streaming, se alimenta
de reparaciones reales y entrega distribuciones de vida residual con
incertidumbre cuantificada.
\end{concepto}

\newpage

%=============================================================================
\section{El problema y la asimetria de costos}
%=============================================================================

El mantenimiento de flota se gestiona hoy con calendarios fijos. Esa
estrategia tiene dos modos de falla simetricos en la descripcion pero
radicalmente asimetricos en el costo:

\begin{itemize}
  \item \textbf{Intervenir antes de tiempo.} Material con vida util restante
    descartado. Costo directo: refaccion, mano de obra, tiempo en taller.
  \item \textbf{Intervenir tarde.} Falla en ruta. Costo multiplicado:
    remolque, downtime no planeado, dano en cadena a componentes adyacentes,
    exposicion regulatoria, carga de trabajo de taller en emergencia.
\end{itemize}

En autotransporte federal mexicano, el costo de una averia en ruta ---
remolque, tiempo muerto, costo de oportunidad del vehiculo y del conductor---
es tipicamente 5 a 20 veces el costo de un servicio preventivo programado.
Esta asimetria es el motor que justifica toda la matematica que sigue.
La metrica correcta de cualquier modelo no es exactitud predictiva;
es \textbf{costo monetario esperado}.

\begin{critico}
\textbf{El error que invalida a la mayoria de los sistemas existentes.}
En una flota viva, la mayoria de las piezas no ha fallado al momento de
observarlas: solo sabemos que su vida $T_i > t_i$ (su edad actual). Esto
es censura por la derecha, y es la regla, no la excepcion. Ignorarla
---por ejemplo, promediar solo las piezas que fallaron--- sesga la
estimacion hacia vidas cortas de forma grosera y produce alertas incorrectas.
La verosimilitud correcta se construye con censura desde el primer dia.
\end{critico}

%=============================================================================
\section{Confiabilidad y vida util remanente}
%=============================================================================

\subsection{La funcion de riesgo: el objeto central}

Sea $T \geq 0$ el tiempo (o km, u horas-motor) hasta la falla de un
componente. Definimos la \textbf{funcion de riesgo instantanea}:

$$h(t) = \lim_{\Delta t \to 0}
  \frac{P(t < T \leq t+\Delta t \mid T > t)}{\Delta t}
  = \frac{f(t)}{R(t)} = -\frac{d}{dt}\ln R(t).$$

Integrar $h$ da toda la distribucion de vida:

$$\boxed{R(t) = \exp\!\left(-\int_0^t h(u)\,du\right).}$$

Especificar $h(t)$ equivale a especificar el modelo completo.
La forma de $h(t)$ tiene fisica adentro: decrece durante mortalidad
infantil, es constante para fallas aleatorias, y \textbf{crece} durante
el desgaste. El regimen creciente es exactamente el que el mantenimiento
predictivo debe anticipar.

\subsection{Weibull: consecuencia del teorema de valores extremos}

Weibull no es una eleccion de conveniencia. Se deriva del
\textbf{teorema de Fisher--Tippett--Gnedenko} para distribuciones de minimos.
Un componente real es una coleccion de $N$ eslabones (granos, asperezas,
volumenes elementales). Falla cuando el mas debil falla:
$T = \min(X_1, \dots, X_N)$.
Si la cola izquierda de cada eslabon es ley de potencia,
$F_X(t) \approx (t/t_0)^\beta$ para $t \to 0^+$:

$$R_T(t) = [1-(t/t_0)^\beta]^N
  \xrightarrow{N \gg 1}
  \exp\!\left[-\left(\frac{t}{\eta}\right)^\beta\right], \quad
  \eta = t_0\,N^{-1/\beta}.$$

El \textbf{parametro de forma} $\beta$ selecciona el regimen:
$\beta < 1$ (mortalidad infantil), $\beta = 1$ (exponencial, sin memoria),
$\beta > 1$ (desgaste; mecanica tipica $\beta \in [1.5,\,4]$).
Consecuencia directa: el efecto de escala $\eta \propto N^{-1/\beta}$
predice que piezas mas grandes son mas debiles, resultado verificado en
materiales y en la literatura tribologica.

\subsection{Vida util remanente: la forma cerrada}

Lo que el operador quiere saber es cuanto le queda a la pieza que
ya sobrevivio hasta la edad $t$. Por integracion por partes:

$$\mathrm{RUL}(t) = \mathbb{E}[T - t \mid T > t]
  = \frac{\displaystyle\int_t^{\infty} R(u)\,du}{R(t)}.$$

Para Weibull con parametros $(\beta,\eta)$, esto tiene forma cerrada:

$$\mathrm{RUL}(t) = \eta\,e^{(t/\eta)^{\beta}}
  \,\Gamma\!\left(1+\tfrac{1}{\beta},\;(t/\eta)^{\beta}\right) - t,$$

donde $\Gamma(a,x) = \int_x^\infty w^{a-1}e^{-w}\,dw$ es la gamma incompleta
superior. Para $\beta > 1$, la RUL decrece estrictamente con la edad: la
pieza vieja tiene menos vida esperada que una nueva. La intervencion se
programa cuando $R(t)$ cae a un nivel objetivo, con un buffer conservador
proporcional a la incertidumbre del estimado.

%=============================================================================
\section{Fisica de falla por componente}
%=============================================================================

La estadistica dice \emph{cuando} falla. La fisica dice \emph{como},
permite construir predictores fisicamente fundamentados, y es lo que
permite generalizar a marcas nuevas sin datos historicos propios.

\subsection{Frenos --- ley de Archard desde el contacto de asperezas}

Dos superficies reales se tocan solo en las cimas de sus rugosidades. Bajo
carga $W$, las asperezas se deforman plasticamente hasta que la presion
de contacto iguala la dureza $H$ del material. El area real de contacto es
$A_r = W/H$, independiente del area aparente (Bowden--Tabor). Al deslizar,
cada junta se rompe en distancia $\sim 2a$; con probabilidad $\kappa$ arranca
una particula hemisferica de volumen $\frac{2}{3}\pi a^3$:

$$\frac{dV}{dL} = \frac{\kappa}{3}\cdot\frac{W}{H}
  \;\equiv\; K\,\frac{W}{H}
  \quad\Longrightarrow\quad
  \boxed{\frac{dh_w}{dt} = k\,p(t)\cdot v(t), \quad k = K/H.}$$

La energia friccional por evento de frenado es
$E_{\rm frenada} \approx \frac{1}{2}m(v_i^2 - v_f^2)$, y el desgaste
acumulado es proporcional a la energia disipada total. Este acumulador ---
construido de IMU (desaceleracion) + GPS (velocidad) + J1939 (masa)---
es el predictor fisico correcto del desgaste de balata. El odometro solo
es un proxy debil.

\subsection{Fatiga estructural --- Miner, Basquin y conteo rainflow}

El dano por carga variable se acumula segun Palmgren--Miner:
$$D = \sum_i \frac{n_i}{N_i} \geq 1 \;\Rightarrow\; \text{falla},
  \qquad N_i \text{ de la curva de Basquin: }
  \sigma_a = \sigma'_f\,(2N_f)^b.$$

La senal continua del IMU $a(t)$ se descompone en ciclos de amplitud
por \textbf{conteo rainflow}: identifica lazos cerrados de histeresis en
la senal, cada uno equivalente a un ciclo de amplitud definida. El dano
acumulado en funcion de los rangos $h_i$ de esos ciclos:
$$D_\beta = \alpha\sum_i h_i^\beta,
  \qquad \beta \approx 3\;\text{(crecimiento de grieta)},
  \quad \beta \approx 5\;\text{(iniciacion)}.$$

El rainflow se calcula en el borde (el IMU no se transmite crudo por LTE).
El resultado es un indice de dano por fatiga por unidad que ningun sistema
basado solo en km puede producir.

\subsection{Rodamientos --- Lundberg--Palmgren como caso especial de Weibull}

La falla por fatiga de contacto se inicia bajo la superficie, donde el
esfuerzo cortante ortogonal $\tau_0$ es maximo. Aplicando el argumento
de eslabon mas debil al volumen esforzado:

$$\ln\frac{1}{S} \propto \frac{\tau_0^c\,N^e\,V}{z_0^h}
  \quad\Longrightarrow\quad
  \boxed{L_{10} = \left(\frac{C}{P}\right)^p,
  \quad p=3\;\text{(bolas)},\quad p=\tfrac{10}{3}\;\text{(rodillos)}.}$$

$C$ es dato del fabricante; $P$ se calcula de las cargas medidas por J1939.
Esta es la misma derivacion de la Seccion~II aplicada al volumen subsuperficial.
La vida $L_{10}$ es exactamente la vida $B_{10}$ de Weibull: la conexion
no es coincidencia sino la misma teoria.

\subsection{Aftertreatment --- el quick win}

J1939 difunde sin pollear: diferencial de presion del DPF (SPN~3251),
nivel de DEF (SPN~1761), derate por NOx (SPN~5246), presion y temperatura
de aceite (SPN~100, 175). La tendencia de $\Delta P$, el ritmo de
regeneraciones forzadas y un codigo pendiente de NOx son senales predictivas
dias antes de que el sistema entre en derate. \textbf{Esta rama es totalmente
alimentable hoy, sin calibracion adicional, y ataca los modos de falla
mas caros del diesel moderno.}

%=============================================================================
\section{Inferencia con datos censurados y covariables}
%=============================================================================

\subsection{La verosimilitud correcta}

Sea $\delta_i = 1$ si la pieza $i$ fallo y $\delta_i = 0$ si sigue viva.
La verosimilitud con censura es:

$$\mathcal{L}(\theta) = \prod_{i=1}^{n}
  h(t_i;\theta)^{\delta_i}\,R(t_i;\theta).$$

Cada observacion aporta su supervivencia $R(t_i)$; las que fallaron
aportan ademas el riesgo $h(t_i)$ en el instante exacto. El log-verosimilitud es
$\ell(\theta) = \sum_i[\delta_i \ln h(t_i;\theta) - H(t_i;\theta)]$.

\subsection{Regresion de riesgos proporcionales de Cox}

Para que la vida dependa de la operacion real, el modelo de
riesgos proporcionales postula:

$$h(t \mid \mathbf{x}) = h_0(t)\,\exp(\boldsymbol{\beta}^{\top}\mathbf{x}),$$

donde $\mathbf{x}$ incluye energia de frenado acumulada, dano de Miner,
severidad de ruta, carga media, marca y conductor. La verosimilitud parcial
de Cox estima $\boldsymbol{\beta}$ sin especificar $h_0(t)$: en cada tiempo
de falla $t_{(j)}$, la probabilidad de que falle la pieza $i$ dado el
conjunto en riesgo $\mathcal{R}_j$ es:

$$P(\text{falla}\;i \mid t_{(j)}) =
  \frac{e^{\boldsymbol{\beta}^\top\mathbf{x}_i}}
       {\sum_{k \in \mathcal{R}_j} e^{\boldsymbol{\beta}^\top\mathbf{x}_k}}.$$

El riesgo base $h_0(t_{(j)})$ se cancela. Se estima el efecto de cada
covariable sobre la vida con minimas suposiciones distribucionales.

%=============================================================================
\section{Procesos de degradacion: la RUL como distribucion completa}
%=============================================================================

Cuando se mide directamente una variable de salud que degrada
($\Delta P_{\rm DPF}$, espesor de balata, capacidad de bateria), la RUL
pasa de ser un estadistico poblacional a una distribucion condicionada al
estado actual del individuo.

\textbf{Proceso de Wiener.} Modelo de degradacion lineal con ruido:
$X(t) = x_0 + \mu t + \sigma B(t)$.
La falla es el primer cruce del umbral $a$. Por el principio de reflexion
del movimiento browniano con deriva:

$$f_{T_a}(t) = \frac{(a-x_0)}{\sigma\sqrt{2\pi t^3}}\,
  \exp\!\left[-\frac{(a-x_0-\mu t)^2}{2\sigma^2 t}\right]
  \;\sim\;
  \mathrm{IG}\!\left(\frac{a-x_0}{\mu},\;\frac{(a-x_0)^2}{\sigma^2}\right).$$

La RUL tiene distribucion \textbf{Gaussiana Inversa}. No es un numero
puntual: es una distribucion completa. Los intervalos de confianza del
pronostico entran directamente en la decision de costo optimo.

La deriva $\mu$ no es un parametro libre: la fija la fisica. En frenos,
$\mu = k\,\langle p\,v \rangle$ (Archard); en el DPF, $\mu$ es la tasa
de acumulacion de hollin. La fisica provee el prior estructural; los datos
ajustan el residuo y la incertidumbre $\sigma$.

\textbf{Proceso Gamma.} Para desgaste estrictamente monotono, el proceso
correcto tiene incrementos no negativos: $\Delta X \sim \mathrm{Gamma}(\alpha\Delta t,\lambda)$.
La confiabilidad en $t$ es la gamma incompleta regularizada evaluada en $a$.

%=============================================================================
\section{Estimacion online en regimen de streaming}
%=============================================================================

Una condicion del mundo real que el diseno debe enfrentar desde el principio:
\textbf{el sistema arranca en un entorno donde la mayoria de las
intervenciones son reactivas.} Las reparaciones no llegan en forma de
datos limpios y etiquetados; llegan como eventos esporadicos, con ordenes
de trabajo de calidad variable, en una flota que mezcla historiales
incompletos con unidades nuevas sin historia alguna. El modelo no puede
esperar a tener datos maduros: debe operar desde el dia uno, actualizar
sus distribuciones conforme lleguen eventos, y mejorar de forma continua
a medida que la flota acumula historia.

La respuesta matematica es el \textbf{filtro bayesiano recursivo},
organizado en capas a distinta cadencia. Ninguna recalcula todo.

\subsection{La ecuacion de Fokker--Planck como motor del filtro}

El estado de degradacion de cada unidad es una densidad de probabilidad
que evoluciona en el tiempo. Para la dinamica de Wiener
$dX = \mu\,dt + \sigma\,dW$, la densidad de creencia $p(x,t)$ obedece
la \textbf{ecuacion de Fokker--Planck}:

$$\frac{\partial p}{\partial t} =
  -\frac{\partial}{\partial x}[\mu\,p]
  + \frac{\sigma^2}{2}\frac{\partial^2 p}{\partial x^2}.$$

El filtro opera en dos pasos por cada observacion:

\begin{enumerate}
  \item \textbf{Prediccion:} propagar $p(x,t)$ segun Fokker--Planck
    hasta el siguiente instante de observacion.
  \item \textbf{Correccion:} actualizar con la observacion $y_t$ por
    Bayes puntual:
    $p(x_t \mid y_{1:t}) \propto p(y_t \mid x_t)\,p(x_t \mid y_{1:t-1})$.
\end{enumerate}

En la practica se implementa como filtro de Kalman (lineal-gaussiano),
UKF (no lineal, momentos exactos de segundo orden), o filtro de particulas
(no parametrico, a costa de costo computacional).

\subsection{Sistemas reparables: el factor de calidad de Kijima}

Una reparacion no produce necesariamente una pieza nueva. El modelo de
\textbf{edad virtual generalizada (GRP) de Kijima} parametriza la eficacia
de cada intervencion:

$$v_{n+1} = q\,(v_n + T_n),$$

donde $v_n$ es la edad virtual antes de la $n$-esima reparacion y $T_n$
el tiempo entre fallas. Con $q = 0$ la reparacion es perfecta
(as-good-as-new); con $q = 1$ es minima (as-bad-as-old); $q > 0$ en la
practica refleja que la calidad de la reparacion es imperfecta.

El sistema \textbf{estima $q$ por taller, por tipo de intervencion y
por componente} a partir de la historia de reincidencia. Esto convierte
cada orden de trabajo en una actualizacion del modelo y produce un
scorecard de calidad de reparacion que ningun sistema de mantenimiento
convencional extrae.

Critico en el regimen de arranque: cuando casi todo es reactivo, la
secuencia de tiempos entre fallas $T_1, T_2, \dots$ es la senal primaria
de aprendizaje. El proceso de renovacion generalizada de Kijima extrae
de esa secuencia tanto la distribucion de vida como la calidad de las
intervenciones, con censura correcta sobre las piezas que aun no vuelven
a fallar.

\subsection{Actualizacion de hiperparametros de flota: SG-MCMC}

Los hiperparametros jerarquicos ($\mu_\beta$, $\tau_\bullet$, etc.) se
actualizan en streaming sin recalcular toda la cadena de Markov. La clave
es que el gradiente estocastico de la log-verosimilitud define una SDE
cuya \textbf{distribucion estacionaria es el posterior}
(Stochastic Gradient Langevin Dynamics):

$$d\theta = \frac{\varepsilon}{2}
  \nabla_\theta \log p(\theta \mid \mathcal{D})\,dt + \sqrt{\varepsilon}\,dW.$$

Minibatches de reparaciones recientes actualizan los hiperparametros sin
acceder a toda la historia. Para flotas no estacionarias (cambio de
flota, nueva normativa, diferente patron de carga estacional) se anade
un factor de olvido exponencial sobre observaciones antiguas.

\subsection{Inferencia amortizada: despliegue instantaneo para unidades nuevas}

Reestimar el modelo para cada unidad nueva con MCMC es prohibitivo en
produccion. La solucion es \textbf{inferencia basada en simulacion (SBI)
/ Neural Posterior Estimation}: entrenar una vez, offline, una red
$q_\psi(\theta \mid x)$ que aprende el mapa datos $\to$ posterior,
con pares $(\theta, x)$ generados por el simulador fisico:

$$q_\psi(\theta \mid x) \approx p(\theta \mid x).$$

En despliegue, la inferencia del posterior de una unidad nueva es un
\textbf{forward pass}, no una corrida de MCMC. El simulador ya existe:
la fisica de falla genera datos sinteticos etiquetados ilimitados.
La red amortizada convierte la inferencia de minutos en milisegundos,
lo cual es necesario para que el sistema escale a miles de unidades.

\medskip

\begin{center}
\small
\begin{tabular}{@{}p{3.4cm}p{4.0cm}p{2.8cm}p{2.4cm}@{}}
\toprule
\small\textbf{Capa} & \small\textbf{Que estima} & \small\textbf{Metodo}
  & \small\textbf{Cadencia} \\
\midrule
\rowcolor{rowalt}
\small Filtrado de estado & \small RUL por unidad & \small KF / UKF / particulas
  & \small Por evento \\
\small Eventos recurrentes & \small Intensidad; $q$ de Kijima & \small GRP / NHPP
  & \small Por reparacion \\
\rowcolor{rowalt}
\small Hiperparametros de flota & \small $(\beta, \eta, \tau_\bullet)$ jerarquicos
  & \small HMC / SG-MCMC & \small Nocturno \\
\small Despliegue & \small Posterior de unidad nueva & \small SBI / forward pass
  & \small Instantaneo \\
\rowcolor{rowalt}
\small Politica & \small Accion optima & \small POMDP / RL & \small Por decision \\
\bottomrule
\end{tabular}
\end{center}

%=============================================================================
\section{Marco jerarquico bayesiano de flota multimarca}
%=============================================================================

Una flota real mezcla marcas, modelos, conductores y rutas. El partial
pooling bayesiano anidado resuelve el problema de datos escasos por
subsegmento sin sacrificar la heterogeneidad real:

$$\ln\eta_i
  = \underbrace{\alpha_{c,k}}_{\text{clase}}
  + \underbrace{\delta_{c,k,m}}_{\text{marca}}
  + \underbrace{\zeta_{c,k,m,j}}_{\text{modelo}}
  + \boldsymbol{\gamma}_c^\top\mathbf{x}_i,$$

$$\alpha_{c,k} \sim \mathcal{N}(\mu_c,\;\tau_{\rm clase}^2), \quad
  \delta \sim \mathcal{N}(0,\;\tau_{\rm marca}^2), \quad
  \zeta \sim \mathcal{N}(0,\;\tau_{\rm modelo}^2).$$

Las $\tau_\bullet$ se estiman de los datos. Si las marcas de una clase se
parecen, el modelo agrupa fuerte; si difieren, desagrupa. No hay eleccion
manual de cuanto pooling aplicar.

Una unidad con menos sensores (OBD vs.\ J1939, o sin CAN) trata la
covariable faltante como variable latente y la marginaliza:
$$p(T_i \mid \mathbf{x}_{i,\mathrm{obs}})
  = \int p(T_i \mid \mathbf{x}_{i,\mathrm{obs}},\,x_{i,\mathrm{mis}})
  \,p(x_{i,\mathrm{mis}} \mid k)\,dx_{i,\mathrm{mis}}.$$

Una unidad con menos informacion produce un posterior predictivo mas ancho.
La incertidumbre se propaga correctamente a la RUL y al umbral de decision.
No hay imputation manual ni descarte de unidades incompletas.

%=============================================================================
\section{Cuantificacion del ahorro}
%=============================================================================

Esta seccion deriva el ahorro esperado en forma cerrada. Los parametros
numericos son estimaciones conservadoras basadas en literatura de
autotransporte; se reemplazan con datos propios de la flota una vez
disponibles.

\subsection{Modelo de costo de referencia: mantenimiento reactivo}

En un regimen puramente reactivo, la tasa de costo de largo plazo
bajo el proceso de renovacion es el teorema de la teoria de renovacion
aplicado al proceso de falla con distribucion Weibull$(\beta, \eta)$:

$$g_{\rm reactivo} = \frac{c_f}{\eta\,\Gamma(1 + 1/\beta)},$$

donde $c_f$ es el costo de una falla no planificada (reparacion + downtime
+ remolque) y $\eta\,\Gamma(1+1/\beta) = \mathrm{MTTF}$ es la vida media.

\subsection{Modelo de costo con mantenimiento preventivo optimo}

Bajo una politica de reemplazo por edad $T$, el costo de largo plazo
por unidad de tiempo es:

$$g(T) = \frac{c_p\,R(T) + c_f\,[1-R(T)]}
  {\displaystyle\int_0^T R(t)\,dt}.$$

La condicion de optimalidad $g'(T^*) = 0$ da:

$$h(T^*)\int_0^{T^*} R(t)\,dt - F(T^*)
  = \frac{c_p}{c_f - c_p}.$$

El lado izquierdo crece con $T$ si y solo si $h$ es creciente ($\beta > 1$),
garantizando que existe un unico minimo global $T^*$. El \textbf{ahorro
unitario por unidad de tiempo} frente a la estrategia reactiva pura es:

$$\Delta g = g_{\rm reactivo} - g(T^*)
  = \frac{c_f}{\eta\,\Gamma(1+1/\beta)} - g(T^*).$$

\subsection{Efecto del error de censura: el costo oculto}

Si el modelo ignora la censura y estima $\hat{\eta}_{\rm sesgado}$
en lugar del verdadero $\eta$, el intervalo de intervencion se desplaza a
$\hat{T}^*$ con costo:

$$g(\hat{T}^*) = g(T^*) + \underbrace{\frac{1}{2}\,g''(T^*)\,(\hat{T}^* - T^*)^2}_{\text{costo del sesgo}}.$$

Para parametros tipicos ($\beta = 2$, fraccion censurada $\sim 70\%$), el
sesgo en $\hat{\eta}$ es del orden del 30--50\% (el modelo ve solo
las piezas que fallaron, que son las de vida corta). El costo del sesgo
es cuadratico en el desplazamiento: un error de 30\% en $\eta$
se traduce en un costo de largo plazo entre 15 y 40\% por encima del optimo.

\subsection{Modelo de ahorro con RUL predictiva (caso general)}

El sistema predice la distribucion de vida residual $p_T(\tau \mid t,\mathbf{x})$
para cada unidad individual, condicionada a su telemetria actual. La politica
optima elige $T^*(t, \mathbf{x})$ que minimiza el costo esperado:

$$g^*(t, \mathbf{x}) = \min_{\tau > 0}\,
  \mathbb{E}_T\!\left[
    \frac{c_p\,\mathbf{1}(T > \tau) + c_f\,\mathbf{1}(T \leq \tau)}
    {\min(T,\tau)}
  \right].$$

La ganancia frente al preventivo optimo de poblacion $g(T^*)$ es:

$$\Delta g_{\rm predictivo}
  = g(T^*) - \mathbb{E}_{\mathbf{x}}\!\left[g^*(t, \mathbf{x})\right],$$

que es estrictamente positiva cuando la variabilidad entre unidades es
no trivial, porque conocer el estado individual permite intervenir mas
tarde en las piezas sanas y mas pronto en las deterioradas.

\subsection{Estimacion numerica con parametros representativos}

Usamos un modelo de balata de freno como caso base, con parametros
conservadores documentados en la literatura de autotransporte:

\begin{center}
\small
\begin{tabular}{@{}lll@{}}
\toprule
\textbf{Parametro} & \textbf{Valor} & \textbf{Fuente} \\
\midrule
$\beta$ (Weibull, balata) & $2.2$ & Tipico de desgaste abrasivo \\
$\eta$ (vida caracteristica) & $80{,}000$ km & Referencia OEM media \\
$c_p$ (preventivo) & \$3{,}500 MXN & Refaccion + mano de obra \\
$c_f$ (correctivo) & \$28{,}000 MXN & $c_p$ + downtime + remolque \\
Flota de referencia & 50 unidades & Flota mediana tipica \\
\bottomrule
\end{tabular}
\end{center}

Con estos valores, el MTTF es:
$$\mathrm{MTTF} = \eta\,\Gamma(1+1/\beta) = 80{,}000 \cdot \Gamma(1.455) \approx 71{,}200\text{ km}.$$

El costo reactivo de largo plazo (por unidad, por km):
$$g_{\rm reactivo} = \frac{28{,}000}{71{,}200} \approx 0.393\text{ MXN/km}.$$

Resolviendo $g'(T^*)=0$ numericamente con los parametros dados:
$T^* \approx 58{,}000$ km, con:
$$g(T^*) = \frac{3{,}500 \cdot R(T^*) + 28{,}000 \cdot [1-R(T^*)]}
  {\displaystyle\int_0^{T^*} R(t)\,dt}
  \approx 0.268\text{ MXN/km}.$$

La reduccion de costo frente a la estrategia reactiva pura es:

\begin{ahorro}
$$\Delta g = 0.393 - 0.268 = \mathbf{0.125\text{ MXN/km por unidad.}}$$

Para una flota de 50 unidades con 60,000 km/ano cada una:

$$\text{Ahorro anual} = 0.125 \times 60{,}000 \times 50
  \approx \mathbf{\$375{,}000\text{ MXN/ano}} \quad (\approx \$19{,}000\text{ USD}).$$

Esto considera solo balatas de eje delantero. Con los demas componentes
(neumáticos, aceite, aftertreatment, batería), la magnitud total es tipicamente
3 a 5 veces mayor, alcanzando \$1.1 a \$1.9 millones MXN por ano
para una flota de 50 unidades.
\end{ahorro}

\subsection{El bono del aftertreatment: un DPF es un caso diferente}

Un filtro de particulas forzado a reparacion correctiva (derate + remolque +
sustitucion del DPF) cuesta entre \$35,000 y \$80,000 MXN por evento,
versus una limpieza preventiva de \$2,500 MXN. Con una tasa de falla de
DPF de $\sim$0.15 eventos/ano por unidad en operacion de ciclo severo, el
ahorro esperado por unidad por ano es:

$$\Delta c_{\rm DPF} = (0.15)(35{,}000 - 2{,}500) = \$4{,}875\text{ MXN/ano.}$$

Para 50 unidades: $\mathbf{\$243{,}750\text{ MXN/ano}}$. Pero el numero mas
relevante es que un DPF correctivo equivale a 14 limpiezas preventivas,
y la distincion entre los dos regimenes es exactamente lo que el sistema
detecta con dias de anticipacion a traves de la tendencia de $\Delta P_{\rm DPF}$.

\subsection{Periodo de convergencia: arranque reactivo}

La transicion del regimen reactivo al predictivo no es instantanea.
El modelo necesita acumular historia de fallas y reparaciones para que los
parametros Weibull converjan. El tiempo de convergencia depende del numero
de eventos observados por componente. Con el estimador de maxima
verosimilitud bajo censura, la varianza del estimado de $\beta$ es:

$$\mathrm{Var}[\hat\beta] \approx \frac{\beta^2}{d},$$

donde $d$ es el numero de fallas observadas (no de unidades). Con $d = 20$
fallas, el error estandar relativo en $\beta$ es $\sim 22\%$; con $d = 50$,
$\sim 14\%$. Para una flota de 50 unidades con tasa de falla de 0.8/ano
por componente, se acumulan 40 fallas en el primer ano: el estimador de
$\beta$ converge a precision util en los primeros 12 a 18 meses.

Durante ese periodo, el modelo opera con el prior jerarquico de clase
(parametros derivados de la literatura de fabricantes), luego va absorbiendo
la historia propia. El costo de ese periodo de convergencia no es cero,
pero tampoco es el costo reactivo completo: el modelo ya alerta sobre
senales directas del bus (DTC, tendencia de $\Delta P$, voltaje de bateria)
que no requieren historia de fallas.

%=============================================================================
\section{Por que este enfoque domina}
%=============================================================================

\subsection{Contra los sistemas basados en reglas OEM}

Los sistemas de reglas (alertar cuando el DTC X aparece, cambiar aceite a
$N$ km) son el punto de partida correcto pero tienen un techo bajo:
no aprenden de la operacion real, no distinguen severidad de ruta,
no cuantifican incertidumbre, y no optimizan el momento de intervencion.
Son un umbral fijo en un problema que tiene estructura continua.
El sistema que describimos los subsume: las reglas OEM son el prior
del modelo antes de ver datos. Una vez que hay historia, el posterior
toma el control.

\subsection{Contra los sistemas puramente de datos (ML sin fisica)}

Los sistemas de ML sin fisica necesitan miles de fallas por componente
y por tipo de vehiculo para generalizar. La razon es que el espacio de
configuraciones (marca, modelo, conductor, ruta, carga) es demasiado grande
para que un modelo sin estructura lo cubra eficientemente. Las dos
consecuencias practicas son: (1) no funcionan en flotas pequeñas o
heterogeneas; (2) no generalizan a marcas nuevas sin reentrenar desde cero.

Nuestro enfoque usa la fisica (Archard, Miner, Lundberg--Palmgren) para
fijar la estructura funcional de la degradacion. El ML estima los residuos
y los parametros de dispersion. Esto reduce la cantidad de datos necesaria
en un orden de magnitud, porque la fisica ya le dice al modelo la forma
correcta de la curva de degradacion.

\subsection{Contra los sistemas de supervivencia sin jerarquia}

Un modelo Weibull plano (sin jerarquia) o un Cox sin partial pooling
sufre el mismo problema: o bien colapsa todo en un modelo global que
ignora la heterogeneidad real entre unidades, o bien ajusta un modelo
separado por subsegmento y se queda sin datos. El partial pooling bayesiano
es la unica solucion que interpola entre estos dos extremos de forma
optima, con la cantidad de agrupamiento determinada por los datos.

\subsection{Contra los sistemas sin cuantificacion de incertidumbre}

Un sistema que reporta ``la balata dura 30 dias mas'' sin intervalo de
confianza es inutilizable para decision de costo optimo. El umbral de
alerta correcto depende de la razon $c_f / c_p$; calcularlo requiere
la distribucion de RUL, no su media. La calibracion de las probabilidades
---que la frecuencia de falla entre unidades con probabilidad predicha
$p$ sea efectivamente $p$--- es tan importante como el error medio de
prediccion, y es lo que el sistema mide explicitamente.

\vspace{0.8cm}

{\color{rule}\rule{\linewidth}{0.4pt}}

\vspace{0.4cm}

\begin{center}
{\small\color{muted}
Junio 2026 \quad|\quad Confidencial \quad|\quad
Para discusion tecnica con equipos de ingenieria}
\end{center}

\end{document}
